% TOP_PROBLEM: {"id": 2307059, "title": "Research Problems in Function Theory — Problem 7.59", "category_id": 8, "category": {"id": 8, "name": "analysis", "display_name": "Analysis", "description": "Limits, continuity, calculus, and function theory.", "slug": "analysis", "order_index": 8, "created_at": "2026-07-31T15:26:25.671Z"}, "status": "open"}
% FIRST_POSTED: null
% ENTRY_KIND: statement_only
% Imported from: batch06.tex; UnsolvedMath ID 2307059
% Compile twice: lualatex 2307059.tex
\documentclass[11pt]{article}
\usepackage{fontspec}
\setmainfont{Latin Modern Roman}
\usepackage{amsmath,amssymb,mathtools,unicode-math}
\setmathfont{Latin Modern Math}
\usepackage{xurl,hyperref}
\hypersetup{hidelinks,pdftitle={UnsolvedMath Problem 2307059}}
\newcommand{\sref}[2]{\hyperlink{source:#1:#2}{[S#2]}}
\newcommand{\cataloguescope}[1]{\paragraph{Mathematical question.}#1}
\begin{document}
% UnsolvedMath ID 2307059; source catalogue code AMR-022-7059
\setcounter{section}{2307058}
\section{Research Problems in Function Theory — Problem 7.59}\label{2307059}
{\small\sffamily UnsolvedMath ID 2307059 · Analysis}
\subsection{Definitions and mathematical statement}

\paragraph{Shared notation.}$\mathbb N=\{1,2,\ldots\}$, and $\mathbb Z,\mathbb Q,\mathbb R,\mathbb C$ denote the integers, rationals, reals and complex numbers. Any other conventions are specified locally.


For $n\ge1$ let $\mathcal E=\mathcal O(\mathbb C^n)$ be all entire functions. For a multi-index $\alpha=(\alpha_1,\ldots,\alpha_n)\in\{0,1,\ldots\}^n$, put $z^\alpha=\prod z_j^{\alpha_j}$ and $D^\alpha=\prod(\partial/\partial z_j)^{\alpha_j}$. If $P(z)=\sum_\alpha p_\alpha z^\alpha$ is a polynomial, define $P(D)=\sum_\alpha p_\alpha D^\alpha$ and $T_{P,Q}f=P(D)(Qf)$. Entire exponential type means $|f(z)|\le A e^{B|z|}$ for some finite positive $A,B$.
\paragraph{Statement recorded in the supplied source.}
For arbitrary complex polynomials $P,Q$ in $n$ variables, does bijectivity of $T_{P,Q}:\mathcal E\to\mathcal E$ imply bijectivity of $T_{Q,P}:\mathcal E\to\mathcal E$? The question concerns all entire functions, not only those of exponential type. The source records that the reversed operator is bijective on the exponential-type class under the hypothesis, but its update gives a negative answer on the full entire class. \sref{2307059}{2}
\subsection{Short English statement}
Historical question: can one interchange the multiplication polynomial and the constant-coefficient differential polynomial while preserving bijectivity on all entire functions?
\subsection{Sources}
\begin{description}
\item[\hypertarget{source:2307059:1}{S1}] ULAM.AI and the UnsolvedMath Contributors, UnsolvedMath: A Curated Collection of Open Mathematics Problems, record 2307059 (AMR-022-7059). CC BY 4.0. \url{https://huggingface.co/datasets/ulamai/UnsolvedMath}.
\item[\hypertarget{source:2307059:2}{S2}] Walter K. Hayman and Edward F. Lingham, Research Problems in Function Theory (2018), Problem 7.59, its complete update and relevant preceding definitions. \url{https://arxiv.org/abs/1809.07200}.
\end{description}
\label{end:2307059}
\end{document}
